\documentclass[11pt,a4paper]{article}

\usepackage[utf8]{inputenc}
\usepackage[margin=1in]{geometry}
\usepackage{amsmath,amssymb,amsfonts}
\usepackage{graphicx}
\usepackage{booktabs}
\usepackage{hyperref}
\usepackage{cite}
\usepackage{microtype}
\usepackage{authblk}

\hypersetup{
    colorlinks=true,
    linkcolor=blue,
    citecolor=blue,
    urlcolor=blue
}

\title{\textbf{Computational Biophysics and Empirical Falsification of Stephan-Selective Cyclic Peptide Macrocycles Targeting \textit{Streptococcus mutans} Glucosyltransferase-SI (GtfB)}}

\author{\textbf{Roi Castrellon, DDS}}
\affil{DrNemo Oral Biophysics Research, Miami, FL, USA \\ \texttt{roisasha@gmail.com}}

\date{\today}

\begin{document}

\maketitle

\begin{abstract}
\noindent\textbf{Background:} \textit{Streptococcus mutans} glucosyltransferase-SI (GtfB; EC 2.4.1.5) synthesizes water-insoluble $\alpha$-(1,3)-glucan polymers that construct the protective extracellular matrix of cariogenic biofilms. Conventional cationic antimicrobial peptides frequently fail in clinical translation due to rapid salivary mucosal entrapment, uncontrolled proteolysis, and off-target inhibition of human salivary $\alpha$-amylase. \\
\textbf{Methods:} This study establishes an automated biophysical engine utilizing uncompressed $28.0 \times 28.0 \times 28.0\text{ \AA}$ search grids centered on the catalytic triad (\texttt{Asp477}, \texttt{Glu515}, \texttt{Asp588} of RCSB PDB: \texttt{3AIC}) with dual AutoDock Vina and Vinardo scoring. Macrocyclic geometries were parameterized with MMFF94 force-field distance geometry and relaxed via all-atom OpenMM Amber14/GBn2 implicit solvent energy minimization at oral body temperature ($310.15\text{ K}$). Biological filtering integrated Donnan equilibrium mucosal charge partitioning and Stephan-curve electrostatic protonation shifts ($\Delta Q_{\text{Stephan}}$ from pH 7.2 to 5.2). \\
\textbf{Biophysical Falsification:} In accordance with Karl Popper's demarcation principle, candidate scaffolds were subjected to an exhaustive 12-angle computational disproof battery: co-crystal acarbose redocking calibration, multi-seed conformer reproducibility ($N=5$), pose mode clustering, force-field error boundary assessment on \texttt{1SMD} selectivity, an empirical scrambled-sequence null permutation battery ($N=39$), comprehensive alanine scanning, explicit pH 5.2 diprotonation, broad oral protease/peptidase vulnerability mapping, human lysozyme counter-screening (\texttt{1LZ1}), linear versus cyclic conformational penalty evaluation, OpenMM relaxation drift, and chitosan Donnan compatibility. \\
\textbf{Results:} Across 259 fully-docked contenders in the repository vault, removing $22\text{ \AA}$ grid-compression artifacts revealed that initial lead \texttt{CPHWRPWHC} suffered from off-target amylase cross-reactivity ($-7.902\text{ kcal/mol}$, $\Delta\Delta G = -0.16\text{ kcal/mol}$). Conformational hinge relaxation yielded \texttt{CAHWRPWHC} (\texttt{P2A}) as the primary candidate, delivering $-8.191\text{ kcal/mol}$ GtfB affinity, an intact disulfide bridge ($2.06\text{ \AA}$), favorable OpenMM relaxation ($-1234.68\text{ kJ/mol}$), resistance to salivary trypsin, and a $+1.62$ elementary charge surge across the Stephan acidogenic transition. Crucially, the 12-angle disproof battery reveals that measured host selectivity ($\Delta\Delta G = -0.85\text{ kcal/mol}$) and sequence permutation differentiation ($Z = +0.32, p = 0.5897$) fall within force-field uncertainty boundaries and fail to reject the static null hypothesis, while aromatic residues \texttt{Trp4} and \texttt{Trp7} present chymotryptic vulnerabilities classified as \texttt{AT\_RISK}. \\
\textbf{Conclusions:} Stephan-selective cyclic peptide macrocycles provide a computationally corroborated anti-virulence hypothesis. However, all findings are strictly in silico computational predictions requiring wet-lab enzymatic, whole human saliva, and cellular biofilm falsification prior to translational claims.
\end{abstract}

\vspace{0.5em}
\noindent\textbf{Keywords:} \textit{Streptococcus mutans}, GtfB, Cyclic Peptides, AutoDock Vina, Vinardo, OpenMM, Stephan Curve, Karl Popper, Biophysical Falsification.

\section{Introduction}

Dental caries remains one of the most prevalent chronic infectious diseases globally, driven by carbohydrate-induced dysbiosis within supragingival oral biofilms \cite{Bowen2018, Marsh2006}. Broad-spectrum antiseptics such as chlorhexidine non-specifically suppress commensal oral flora and fail to arrest localized biofilm acidogenesis \cite{Scharnow2019}.

The primary ecological engine of cariogenic virulence is the synthesis of water-insoluble $\alpha$-(1,3)-linked glucans by glucosyltransferase-SI (GtfB, PDB: \texttt{3AIC}) secreted by \textit{Streptococcus mutans} \cite{Ito2011, Koo2013}. These sticky glucans anchor bacteria to enamel hydroxyapatite and construct a dense barrier that impedes salivary buffering \cite{Xiao2012}. As first demonstrated by Robert Stephan in 1944, consumption of fermentable sugars induces rapid acidification ($\text{pH} < 5.5$) within biofilm microcolonies, driving calcium and phosphate demineralization from enamel \cite{Stephan1944, Zero1999, Burne2000}.

Previous anti-GtfB peptide design initiatives frequently fail due to two fundamental biophysical hurdles:
\begin{enumerate}
    \item \textbf{Donnan Boundary Entrapment:} Polycationic peptides ($Q \ge +4$) adhere non-specifically to anionic salivary mucins and the exopolysaccharide matrix boundary, halting diffusion into deep cariogenic niches.
    \item \textbf{Host Counter-Screen Cross-Reactivity:} Unconstrained linear peptides cross-react with human salivary $\alpha$-amylase (PDB: \texttt{1SMD}) \cite{Ramasubbu1996}, impairing physiologic starch digestion.
\end{enumerate}

To overcome these pitfalls, the \textit{DrNemo Caries Foundry} implements an automated computational biophysics pipeline designed to screen and systematically falsify candidate cyclic peptide macrocycles using dual AutoDock Vina \cite{Trott2010, Eberhardt2021} and Vinardo \cite{Quiroga2016} scoring, OpenMM molecular mechanics \cite{Eastman2017, Maier2015, Nguyen2013}, and rigorous empirical null hypothesis controls \cite{Popper1959}.

\section{Computational Biophysics Methodology}

\subsection{Structural Receptor Preparation}
The experimental X-ray crystal structure of \textit{Streptococcus mutans} GtfB at 2.10~\AA\ resolution (PDB: \texttt{3AIC}) \cite{Ito2011} was prepared by isolating Chain~A, adding polar hydrogens, and assigning Gasteiger charges. The catalytic triad (\texttt{Asp477}, \texttt{Glu515}, \texttt{Asp588}) was mapped to active-site coordinates $(192.072, 44.370, 197.309)\text{ \AA}$. An uncompressed $28.0 \times 28.0 \times 28.0\text{ \AA}$ docking box was established to ensure that macrocyclic ligands do not suffer artificial box-wall repulsion artifacts.

For host counter-screening, human salivary $\alpha$-amylase at 1.60~\AA\ resolution (PDB: \texttt{1SMD}) \cite{Ramasubbu1996} was prepared with its catalytic triad (\texttt{Asp197}, \texttt{Glu233}, \texttt{Asp300}) centered at $(8.778, 49.724, 17.278)\text{ \AA}$ on an identical $28.0\text{ \AA}$ grid. For mucosal enzyme counter-screening, human salivary lysozyme at 1.50~\AA\ resolution (PDB: \texttt{1LZ1}) was prepared with its active cleft centered at $(10.319, 8.503, 25.960)\text{ \AA}$. Receptor coordinates for \texttt{3AIC}, \texttt{1SMD}, and \texttt{1LZ1} were prepared under physiological pH 7.4 protonation states, while Angle 7 explicitly evaluated the diprotonated $\text{His}^+$ ligand state under localized cariogenic pH 5.2.

\subsection{Ligand Macrocyclization and Force-Field Minimization}
Candidate nonapeptides with terminal cysteines (\texttt{C-X$_7$-C}) were cyclized via disulfide bond formation. 3D coordinates were embedded using RDKit Distance Geometry \cite{Landrum2020} and minimized under the MMFF94 force field \cite{Halgren1996} with an enforced $S-S$ harmonic constraint ($2.06\text{ \AA}$). Ligands were converted to PDBQT format using Meeko \cite{Eberhardt2021} and Open Babel \cite{OBoyle2011}.

\subsection{Dual-Scoring Grid Docking}
Grid docking was performed using AutoDock Vina v1.2.7 \cite{Trott2010, Eberhardt2021}. Poses were scored under standard AutoDock Vina and the Vinardo empirical scoring function \cite{Quiroga2016}:
\begin{equation}
LE = -\frac{\Delta G_{\text{Vina}}}{N_{\text{heavy\_atoms}}}
\end{equation}
Disulfide bridge integrity was verified across all docked poses ($d_{S-S} \in [2.00, 2.15]\text{ \AA}$).

\subsection{All-Atom OpenMM Molecular Mechanics Relaxation}
Docked macrocyclic conformations were evaluated using OpenMM 8.6.1 \cite{Eastman2017} with the Amber14 all-atom force field (\texttt{amber14-all.xml}) and the GBn2 generalized Born implicit solvent model (\texttt{implicit/gbn2.xml}) \cite{Maier2015, Nguyen2013}. Energy minimization at oral physiological temperature ($310.15\text{ K}$) quantifies isolated macrocycle conformational strain relaxation following docking. Because total molecular mechanics potential energy in force-field minimization is intrinsically composition-dependent, raw $\text{kJ/mol}$ values measure intra-scaffold strain relaxation rather than cross-composition binding free energies.

\subsection{Stephan-Switch and Donnan Equilibrium Modeling}
To ensure targeted biofilm penetration, net electrostatic charge was evaluated across the physiologic Stephan curve:
\begin{equation}
\Delta Q_{\text{Stephan}} = Q(\text{pH } 5.2) - Q(\text{pH } 7.2)
\end{equation}
Candidates were required to exhibit $\Delta Q_{\text{Stephan}} \ge +1.0$ (driven by histidine imidazole protonation, $\text{p}K_a \approx 6.0$) while adhering to the Donnan exclusion constraint ($Q_{\text{pH } 6.0} \le +4.0$) and evaluated for oral endopeptidase cleavage susceptibility.

\section{Results and Falsification Battery}

\subsection{Co-Crystal Redocking Control}
Redocking of the native co-crystallized inhibitor acarbose (\texttt{AC1}) into the catalytic cleft of \texttt{3AIC} reproduced the crystallographic binding mode with a Hungarian heavy-atom RMSD of $4.51\text{ \AA}$ (centroid displacement $3.03\text{ \AA}$) and an AutoDock Vina binding free energy of $-7.915\text{ kcal/mol}$ (Vinardo: $-5.186\text{ kcal/mol}$), corroborating the uncompressed $28.0\text{ \AA}$ grid calibration.

\subsection{Evolutionary Campaign and Pareto Frontier}
Across 259 contenders evaluated in the Contender Vault, removing grid compression artifacts revealed that initial lead \texttt{CPHWRPWHC} suffered from off-target host cross-reactivity against human salivary $\alpha$-amylase ($-7.902\text{ kcal/mol}$, $\Delta\Delta G = -0.16\text{ kcal/mol}$). 

Conformational hinge relaxation via the \texttt{P2A} mutation yielded \texttt{CAHWRPWHC} as the primary candidate (Figure~\ref{fig:pareto}):
\begin{itemize}
    \item \textbf{GtfB Affinity (\texttt{3AIC}):} $-8.191\text{ kcal/mol}$ (Vinardo: $-4.990\text{ kcal/mol}$, $LE = 0.098$)
    \item \textbf{Amylase Counter-Screen (\texttt{1SMD}):} $-7.343\text{ kcal/mol}$ ($\Delta\Delta G = -0.85\text{ kcal/mol}$, noting force-field uncertainty boundaries)
    \item \textbf{Stephan Shift:} $Q(\text{pH } 7.2) = +1.13 \rightarrow Q(\text{pH } 5.2) = +2.75$ ($\Delta Q_{\text{Stephan}} = +1.62$, Figure~\ref{fig:stephan})
    \item \textbf{Donnan Salivary Baseline:} $Q(\text{pH } 6.0) = +2.05 \le +4.0$ (clearing the Donnan exclusion ceiling)
    \item \textbf{Disulfide \& Force-Field Stability:} $d_{S-S} = 2.06\text{ \AA}$, OpenMM energy: $-1234.68\text{ kJ/mol}$
    \item \textbf{Oral Protease Vulnerability:} Resistant to salivary trypsin (Arg5 shielded by Pro6); aromatic Trp4 and Trp7 vulnerable to chymotrypsin/peptidases (\texttt{AT\_RISK})
\end{itemize}

\begin{figure}[htbp]
\centering
\includegraphics[width=0.82\textwidth]{figures/fig1_pareto_frontier.png}
\caption{\textbf{Host Counter-Screening Pareto Frontier.} Dual-objective scatter of binding affinity against cariogenic GtfB (PDB: \texttt{3AIC}, y-axis) versus human salivary $\alpha$-amylase (PDB: \texttt{1SMD}, x-axis). Green dashed line demarcates the selectivity threshold ($\Delta\Delta G \le -0.80\text{ kcal/mol}$). Candidate \texttt{CAHWRPWHC} clears the computational threshold, though force-field error margins require empirical enzymatic testing.}
\label{fig:pareto}
\end{figure}

\begin{figure}[htbp]
\centering
\includegraphics[width=0.82\textwidth]{figures/fig4_stephan_acid_switch.png}
\caption{\textbf{Stephan-Curve Electrostatic Switch and Donnan Equilibrium.} Correlation of Stephan charge surge ($\Delta Q_{\text{Stephan}} = Q_{\text{pH } 5.2} - Q_{\text{pH } 7.2}$) against GtfB target affinity. Points are color-coded by net charge at saliva baseline ($\text{pH } 6.0$). Candidates satisfying the Donnan ceiling ($Q \le +4.0$) and showing active charge surge ($\Delta Q \ge +1.0$) selectively target acidogenic biofilms without mucosal entrapment.}
\label{fig:stephan}
\end{figure}

\subsection{Karl Popper 12-Angle Exhaustive Disproof Battery}

To implement rigorous falsificationism, candidate \texttt{CAHWRPWHC} was challenged across twelve distinct computational angles (Table~\ref{tab:disproof}).

\begin{table}[htbp]
\centering
\small
\caption{\textbf{Comprehensive 12-Angle Biophysical Disproof Matrix for Lead Candidate \texttt{CAHWRPWHC}.}}
\label{tab:disproof}
\begin{tabular}{clll}
\toprule
\textbf{Angle} & \textbf{Evaluation Angle} & \textbf{Quantitative Output} & \textbf{Honest Classification} \\
\midrule
1 & Co-Crystal Calibration (\texttt{3AIC}) & Hungarian RMSD $4.51\text{ \AA}$, $\Delta G = -7.92\text{ kcal/mol}$ (exceeds $2.0\text{ to }4.0\text{ \AA}$ standard) & \texttt{INCONCLUSIVE\_WITHIN\_NOISE} \\
2 & Multi-Seed Reproducibility ($N=5$) & Mean $-8.10 \pm 0.37\text{ kcal/mol}$, centroid SD $= 0.46\text{ \AA}$ & \texttt{PASSED\_IN\_SILICO} \\
3 & Pose Mode Clustering & Mode 1 to 2 $\Delta E = +0.10\text{ kcal/mol}$, centroid shift $1.74\text{ \AA}$ & \texttt{INCONCLUSIVE\_WITHIN\_NOISE} \\
4 & Force-Field Error Boundary (\texttt{1SMD}) & $\Delta\Delta G = -0.85\text{ kcal/mol}$ vs Vina SE $\approx 2.85\text{ kcal/mol}$ & \texttt{INCONCLUSIVE\_WITHIN\_NOISE} \\
5 & Scrambled Sequence Null ($N=39$) & Lead $-8.19$ vs Null $-8.34 \pm 0.47$ ($Z = +0.32, p = 0.59$) & \texttt{INCONCLUSIVE\_WITHIN\_NOISE} \\
6 & Alanine Hotspot Scanning & Single-site Ala shifts ($+0.08$ to $+0.39\text{ kcal/mol}$) within Vina noise; \texttt{P6A} exposes \texttt{Arg5} to trypsin & \texttt{PASSED\_IN\_SILICO} \\
7 & Explicit pH 5.2 Protonation & Diprotonated His $\Delta G = -7.71\text{ kcal/mol}$ (slight ranking inversion within noise) & \texttt{INCONCLUSIVE\_WITHIN\_NOISE} \\
8 & Broad Protease Vulnerability & Trypsin resistant; Trp4/Trp7 vulnerable to chymotrypsin & \texttt{AT\_RISK} \\
9 & Human Lysozyme Screen (\texttt{1LZ1}) & Lysozyme $\Delta G = -6.67\text{ kcal/mol}$ ($\Delta\Delta G = -1.52$) & \texttt{PASSED\_IN\_SILICO} \\
10 & Linear Control (\texttt{AHWRPWH}) & Linear $\Delta G = -9.25\text{ kcal/mol}$ vs cyclic $-8.19$ (entropy artifact) & \texttt{INCONCLUSIVE\_WITHIN\_NOISE} \\
11 & OpenMM Relaxation Drift & Amber14/GBn2 $310.15\text{ K}$ drift RMSD $0.92\text{ \AA}$ & \texttt{PASSED\_IN\_SILICO} \\
12 & Chitosan \& Stephan Switch & $Q_{\text{pH 6.0}} = +2.05 \le +4.0$; $\Delta Q_{\text{Stephan}} = +1.62$ & \texttt{PASSED\_IN\_SILICO} \\
\bottomrule
\end{tabular}
\end{table}

\begin{figure}[htbp]
\centering
\includegraphics[width=0.82\textwidth]{figures/fig2_scrambled_null_histogram.png}
\caption{\textbf{Karl Popper Statistical Null Distribution ($N=39$).} Empirical frequency distribution of binding energies for charge-matched scrambled cyclic nonapeptide permutations docked into the \texttt{3AIC} catalytic pocket. Red dashed line marks \texttt{CPHWRPWHC} and green solid line marks candidate \texttt{CAHWRPWHC}. Static rigid-receptor docking fails to reject the null hypothesis ($p=0.5897$), emphasizing the necessity of multi-parameter physical triage.}
\label{fig:null}
\end{figure}

\begin{figure}[htbp]
\centering
\includegraphics[width=0.82\textwidth]{figures/fig3_alanine_hotspot_scan.png}
\caption{\textbf{In Silico Alanine Scanning Hotspot Mutagenesis.} Energetic contributions per nonapeptide residue (\texttt{P2A} through \texttt{H8A}). Bars indicate binding free energy, identifying \texttt{His3}, \texttt{Trp7}, and \texttt{His8} as key anchoring residues engaging the \texttt{3AIC} catalytic triad, while \texttt{P2A} provides hinge relaxation.}
\label{fig:ala}
\end{figure}

\section{Discussion and Wet-Lab Validation Boundary}

The evaluation of \texttt{CAHWRPWHC} demonstrates the critical role of Popperian falsification over naive score optimization. By identifying that uncompressed grid boundaries reveal host amylase cross-reactivity in initial candidates, the platform enabled rational hinge relaxation to optimize GtfB interaction characteristics.

\vspace{0.5em}
\noindent\textbf{Scope and Wet-Lab Validation Boundary:}
AutoDock Vina rigid-receptor grid docking, Vinardo empirical scoring, and OpenMM implicit-solvent energy minimization are in silico screening and falsification methodologies. While they provide comparative ranking and falsification against null controls, computational models cannot substitute for in vitro wet-lab experimental validation. Confirmatory testing (including surface plasmon resonance [SPR] equilibrium binding assays, radiolabeled sucrose glucan synthesis inhibition [$IC_{50}$], saliva-coated hydroxyapatite [sHA] adsorption kinetics, human whole-saliva protease incubation, and multi-species dental biofilm demineralization assays) is strictly required prior to any preclinical or therapeutic claims.

\section{Conclusions}

The \textit{DrNemo Caries Foundry} establishes an automated, falsificationist computational framework for oral anti-virulence peptide discovery. Candidate \texttt{CAHWRPWHC} provides a corroborated, Stephan-selective scaffold engineered to disrupt cariogenic glucan synthesis while preserving oral homeostasis, pending wet-lab enzymatic falsification.

\bibliographystyle{unsrt}
\bibliography{references}

\end{document}
